% !TeX root = ./teoremiprincipali.tex
\documentclass{article}
\usepackage[utf8]{inputenc}
\usepackage{mathtools}
\usepackage{amssymb}
\usepackage{centernot}
\usepackage{bm}
\usepackage{fullpage}
\usepackage{multicol}

\begin{document}

\section{Teorema di Bolzano-Weierstrass}

\subsection{Definizioni preliminari}
Nessuna.

\begin{multicols}{2}
    \subsection{Ipotesi}
    Successione \(a_n\) \textbf{limitata} (superiormente e inferiormente).
\columnbreak
    \subsection{Tesi}
    Esistono \textsc{infinite sottosuccessioni} (relative alla successione iniziale) convergenti.
\end{multicols}

\subsection{Dimostrazione}

Chiamiamo \(A_0\) l'insieme che contiene tutti i punti della successione.\\
Eseguiamo la seguente procedura per \(k = 0\).
\begin{enumerate}
    \item Chiamiamo \([\alpha_k, \beta_k]\) i due estremi dell'intervallo della successione.\\
          Prendiamo il punto medio tra i due estremi, e chiamiamolo \(\gamma_k\).\\
          Osserviamo che \(\alpha_k \leq \gamma_k \leq \beta_k\), e che la dimensione \(d_k\) dell'intervallo \([\alpha_k, \gamma_k] = [\gamma_k, \beta_k]\) è la metà di \([\alpha_k, \beta_k]\).
    \item Creiamo due insiemi di punti della successione: uno con i punti tra \([\alpha_k, \gamma_k]\) e uno con i punti tra \(]\gamma_k, \beta_k]\).
    \item Almeno uno dei due insiemi ha un numero infinito di punti: prendiamolo, e chiamiamolo \(A_{k+1}\).
\end{enumerate}
Possiamo ripetere questa procedura un numero infinito di volte: possiamo notare che le dimensioni dell'intervallo \(d_k = (\frac{d_0}{2^k}) \to 0\); dato che \(A_k\) contiene infiniti punti, possiamo creare una sottosuccessione che includa solo punti contenuti in \(A_k\).\\
Essa sarà convergente per il teorema dei carabinieri ad un valore \(L\) tale che \(\alpha_0 \leq \dots \leq \alpha_k \leq L \leq \beta_k \leq \dots \leq \beta_0\).

\newpage

\section{Polinomio di Taylor con resto di Peano}

\subsection{Definizioni preliminari}
\[P_{n, x_0}(x) = (\sum^{n}_{m = 0} \frac{f^{(m)}(x_0) * (x - x_0)^m}{m!})\]

\begin{multicols}{2}
    \subsection{Ipotesi}
    Funzione \(f(x) :\ ]a, b[ \to \mathbb{R}\), \textbf{derivabile} \(n\) volte in \(x_0\) e \(n - 1\) volte in \(]a, b[\).\\
    Punto \(x_0 \in\ ]a, b[\).
\columnbreak
    \subsection{Tesi}
    La funzione \(f(x)\) è \textsc{approssimabile} nel punto \(x_0\) con il polinomio \(P_{n, x_0}(x) + o(x - x_0)^n\) di grado \(n\).
\end{multicols}

\subsection{Dimostrazione}

Notiamo che \(P^{n}_{n, x_0}(x_0) = f^{(n)}(x_0)\).\\
Proviamo a calcolare il seguente limite, che ci sarà utile nel prossimo passaggio:
\[\lim_{x \to x_0} \frac{f(x) - P_{n - 1, x_0}(x)}{(x - x_0)^n} =^{DH} \lim_{x \to x_0} \frac{f^{(1)}(x) - P^{(1)}_{n - 1, x_0}(x)}{n * (x - x_0)^{n - 1}} =^{\infty DH} \lim_{x \to x_0} \frac{f^{(n - 1)}(x) - P^{(n - 1)}_{n - 1, x_0}(x)}{n! * (x - x_0)^{1}} =\]
\[= \lim_{x \to x_0} \frac{f^{(n - 1)}(x) - f^{(n - 1)}(x_0)}{n! * (x - x_0)}\]
Ora siamo pronti a calcolare il limite con \(n\) invece che \(n - 1\):
\[\lim_{x \to x_0} \frac{f(x) - P_{n, x_0}(x)}{(x - x_0)^n}\]
Estraiamo un termine dal polinomio:
\[\lim_{x \to x_0} \frac{f(x) - P_{n - 1, x_0}(x) - \frac{f^{(n)}(x_0) * (x - x_0){n!}}{x - x_0}}{(x - x_0)^n}\]
Raccogliamo termini in modo da formare il limite precedente:
\[\lim_{x \to x_0} ( \frac{f(x) - P_{n - 1, x_0}(x)}{(x - x_0)^n} - \frac{\frac{f^{(n)}(x_0) * (x - x_0)^n}{n!}}{(x - x_0)^n} ) \]
Facciamo uscire dal limite le costanti:
\[- \frac{f^{(n)}(x_0)}{n!} + \lim_{x \to x_0} \frac{f(x) - P_{n - 1, x_0}(x)}{(x - x_0)^n}\]
Per il limite precedente:
\[- \frac{f^{(n)}(x_0)}{n!} + \lim_{x \to x_0} \frac{f^{(n - 1)}(x) - f^{(n - 1)}(x_0)}{n! * (x - x_0)}\]
Raccogliamo \(\frac{1}{n!}\):
\[\frac{1}{n!} (- f^{(n)}(x_0) + \lim_{x \to x_0} \frac{f^{(n - 1)}(x) - f^{(n - 1)}(x_0)}{(x - x_0)})\]
Abbiamo ottenuto un rapporto incrementale, il che significa che:
\[\frac{1}{n!} (- f^{(n)}(x_0) + f^{(n)}(x_0)) = 0\]

\newpage

\section{Teorema di esistenza degli zeri}

\subsection{Definizioni preliminari}
Nessuna.

\begin{multicols}{2}
    \subsection{Ipotesi}
    Funzione \(f(x)\) \textbf{continua} in \([a_0, b_0]\).
\columnbreak
    \subsection{Tesi}
    Esiste \textsc{almeno un punto} in cui \(f(x) = 0\).
\end{multicols}

\subsection{Dimostrazione}

Notiamo che \(f(a_0) * f(b_0) \leq 0\) (ovvero è negativa).\\
Definiamo la seguente procedura:
\begin{enumerate}
    \item Bisezioniamo l'intervallo \([a_n, b_n]\) in \([a_n, z_n]\) e \([z_n, b_n]\).
    \item Almeno uno dei due intervalli è tale che \(f(inizio) * f(fine) \leq 0\) (negativo).
    \item Prendiamo un intervallo per il quale il prodotto precedente è negativo, e chiamiamolo \([a_{n+1}, b_{n+1}]\).
\end{enumerate}
Ripetendo infinite volte la procedura, partendo dall'intervallo \([a_0, b_0]\), otterremo un intervallo sempre più "verticalmente stretto" \([a_n, b_n]\).\\
Possiamo notare che \(a_0 \leq a_n \leq b_n \leq b_0\), e che entrambe le successioni tendono allo stesso numero \(a_n \to x\) e \(b_n \to x\).\\
Possiamo considerare quindi quel numero come punto singolo.\\
Calcoliamo nuovamente \(f(a_n) * f(b_n)\): sappiamo che risulta essere \(\leq 0\), ma possiamo sostituire il limite: \(f(x) * f(x) \leq 0\).\\
Dunque, abbiamo che \(f(x)^2 \leq 0\), e quindi che \(f(x) = 0\).

\end{document}
